\documentclass[12pt,a4paper]{article}
%%%%%%%%%%%%%%%%%%%%%%%%%%%%%%%%%%%%%%%%%%%%
%                 PACKAGES
%%%%%%%%%%%%%%%%%%%%%%%%%%%%%%%%%%%%%%%%%%%%

% LANGUAGE AND ENCODING
\usepackage[utf8]{inputenc}      % Direct UTF-8 character input
\usepackage[english]{babel}      % Language-specific hyphenation
\usepackage{xstring}

% CORE MATHEMATICS
\usepackage{amsmath}             % AMS math environments (align, gather, etc.)
\usepackage{amssymb}             % AMS math symbols
\usepackage{amsthm}              % Theorem environments
\usepackage{amsfonts}            % AMS fonts
\usepackage{mathtools}           % Extensions to amsmath
\usepackage{esint}               % Extended integral symbols
\usepackage{cancel}              % Diagonal cancellation in math
\usepackage{tcolorbox}           % Good boxes
\tcbuselibrary{most}

\usepackage{lmodern}
\usepackage{bm}                  % Bold math symbols (Greek, etc.)
\usepackage{mathrsfs}            % Ralph Smith's script font
\usepackage{tensor}              % Tensor index notation

% UNITS AND QUANTITIES
\usepackage{siunitx}             % Professional number and unit formatting
\NewCommandCopy\uqty\qty         % Save siunitx's \qty as \uqty
\usepackage{silence}
\WarningFilter{siunitx}{Detected the "physics" package}
\ExplSyntaxOn
\msg_redirect_name:nnn { siunitx } { physics-pkg } { none }
\ExplSyntaxOff

% TABLES AND BOXES
\usepackage{tabulary}            % Tables with automatic width adjustment
\usepackage{mdframed}            % Framed boxes with backgrounds
% \usepackage{chemformula}       % Alternative chemistry notation (commented)

% MATH TOOLS
\usepackage{centernot}           % Centered negation symbol
\usepackage{dcolumn}             % Decimal alignment in tables
\usepackage[shortlabels]{enumitem} % Customizable lists

% GRAPHICS AND FIGURES
\usepackage{graphicx}            % Include images
\usepackage{float}               % Force figure placement with [H]
\usepackage{tikz}                % Powerful diagram drawing
\usetikzlibrary{positioning,arrows,trees,3d,angles} % TikZ libraries
\usepackage{pgfplots}            % Scientific plots based on tikz
\pgfplotsset{compat=1.18}        % PGFPlots compatibility mode
\usepackage{subcaption}          % Subfigures (a), (b), etc.
\usepackage{caption}             % Customize caption formatting
\usepackage{booktabs}            % Professional tables with \toprule

% MULTI-COLUMN AND ARRAYS
\usepackage{multicol}            % Multi-column text layout
\usepackage{array}               % Enhanced table column types

% CHEMISTRY
\usepackage{chemfig}             % Draw molecular structures
\usepackage[version=4]{mhchem}   % Chemical equations in text (\ce{H2O})

% CIRCUITS
\usepackage{circuitikz}          % Circuit diagrams with tikz

% TEXT FORMATTING
\usepackage{blindtext}           % Lorem ipsum placeholder text
\usepackage{titlesec}            % Customize section headings
\usepackage{fancyhdr}            % Custom headers and footers
\usepackage{csquotes}            % Context-sensitive quotes

% HYPERLINKS
\usepackage[colorlinks=true, allcolors=blue]{hyperref} % Clickable links

% PAGE GEOMETRY
\usepackage[top=1.0in,bottom=1.0in,left=1.0in,right=1.0in,marginparwidth=1.75cm]{geometry}

% CODE LISTINGS
\usepackage{xcolor}              % Color definitions
\usepackage{listings}            % Code syntax highlighting
\usepackage{xparse}              % Robust command definitions
\usepackage{braket}              % Dirac bra-ket notation

% PHYSICS PACKAGES (LOADED LAST)
\usepackage[italicdiff]{physics} % Original physics package (d in italics)
\usepackage{physics2}            % Modern physics package
\usephysicsmodule{ab}            % Auto-braces module
\usephysicsmodule{ab.legacy}     % Legacy \abs, \norm, \eval
\usephysicsmodule{nabla.legacy}  % \grad, \div, \curl operators
\usephysicsmodule{op.legacy}     % Operators like \tr, \Re, \Im
\usephysicsmodule{diagmat, xmat} % Diagonal and special matrices
\usephysicsmodule{braket, ab.braket} % Dirac notation
\let\lpc\laplacian               % Alias for Laplacian

%%%%%%%%%%%%%%%%%%%%%%%%%%%%%%%%%%%%%%%%%%%%
%             NEW COMMANDS
%%%%%%%%%%%%%%%%%%%%%%%%%%%%%%%%%%%%%%%%%%%%

%%% MATH OPERATORS
\DeclareMathOperator{\sinc}{sinc}
\DeclareMathOperator{\sgn}{sgn}

%%% INVERSE TRIGONOMETRIC OPERATORS
\renewcommand{\asin}{\sin^{-1}}
\renewcommand{\acos}{\cos^{-1}}
\renewcommand{\atan}{\tan^{-1}}
\renewcommand{\acsc}{\csc^{-1}}
\renewcommand{\asec}{\sec^{-1}}
\renewcommand{\acot}{\cot^{-1}}

\newcommand{\asinh}{\sinh^{-1}}
\newcommand{\acosh}{\cosh^{-1}}
\newcommand{\atanh}{\tanh^{-1}}
\newcommand{\acsch}{\csch^{-1}}
\newcommand{\asech}{\sech^{-1}}
\newcommand{\acoth}{\coth^{-1}}
% \DeclareMathOperator{\arcsin}{arcsin}
% \DeclareMathOperator{\arccos}{arccos}
% \DeclareMathOperator{\arctan}{arctan}
% \DeclareMathOperator{\arccsc}{arccsc}
% \DeclareMathOperator{\arcsec}{arcsec}
% \DeclareMathOperator{\arccot}{arccot}
\DeclareMathOperator{\arsinh}{arsinh}
\DeclareMathOperator{\arcosh}{arcosh}
\DeclareMathOperator{\artanh}{artanh}
\DeclareMathOperator{\arcsch}{arcsch}
\DeclareMathOperator{\arsech}{arsech}
\DeclareMathOperator{\arcoth}{arcoth}

%%% NUMBER SETS
\newcommand{\R}{\mathbb{R}}          % Real numbers
\newcommand{\Z}{\mathbb{Z}}          % Integers
\newcommand{\Q}{\mathbb{Q}}          % Rational numbers
\newcommand{\C}{\mathbb{C}}          % Complex numbers
\newcommand{\N}{\mathbb{N}}          % Natural numbers
\newcommand{\K}{\mathbb{K}}          % Field R or C
\newcommand{\bb}[1]{\mathbb{#1}}     % Generic blackboard bold

%%% PARTIAL DIFFERENTIALS
\newcommand{\p}{\partial}            % Partial symbol

%%% VECTOR CALCULUS
\newcommand{\krondelta}[1]{\delta_{#1}}
\newcommand{\levicivita}[1]{\varepsilon_{#1}}
\newcommand{\ihat}{\vu*{\imath}}
\newcommand{\jhat}{\vu*{\jmath}}
\newcommand{\khat}{\vu*{k}}
\newcommand{\rhohat}{\vu*{\rho}}
\newcommand{\phihat}{\vu*{\phi}}
\newcommand{\thetahat}{\vu*{\theta}}
\newcommand{\rhat}{\vu*{r}}
\newcommand{\inv}{^{-1}}
\newcommand{\vell}{\boldsymbol{\ell}}

%%% DIFFERENTIALS AND TRANSFORMS
\newcommand{\laplace}[1]{\mathscr{L}\qty{#1}}
\newcommand{\invlaplace}[1]{\mathscr{L}^{-1}\qty{#1}}
\newcommand{\unitstep}[1]{\mathscr{U}\qty{#1}}
\newcommand{\intfty}{\int_{-\infty}^{+\infty}}

%%% PARENTHESES
\newcommand{\floor}[1]{\lfloor #1 \rfloor}
\newcommand{\ceiling}[1]{\lceil #1 \rceil}
\newcommand{\round}[1]{\lfloor #1 \rceil}
\newcommand{\nvec}[3]{\left\langle #1 , #2 , #3 \right\rangle}
\newcommand{\tline}[1]{\hrule \vspace{#1} \hrule}
\newcommand{\blank}[1]{\hspace*{#1}}

%%% TRANSFORM OPERATORS
\newcommand{\Laplace}{\mathscr{L}}
\newcommand{\Fourier}{\mathscr{F}}

%%% LOGICAL OPERATORS
\newcommand{\union}{\cup}
\newcommand{\intersect}{\cap}
\newcommand{\defined}{:=}

%%% COMPLEX ANALYSIS
\newcommand{\Arg}{\operatorname{Arg}}
\newcommand{\Hol}{\operatorname{Hol}}
\newcommand{\Log}{\operatorname{Log}}

%%% LINEAR ALGEBRA
\newcommand{\Span}[1]{\operatorname{span}\qty(#1)}
\newcommand{\Rank}[1]{\operatorname{rank}\qty(#1)}
\newcommand{\nullity}[1]{\operatorname{nullity}\qty(#1)}
\DeclareMathOperator{\image}{im}

%%% TOPOLOGY
\newcommand{\Id}{\operatorname{Id}}
\newcommand{\Image}[1]{\operatorname{Im}\qty(#1)}
\newcommand{\comp}{^{c}}
\newcommand{\Prob}[1]{\operatorname{Pr}\qty(#1)}
\newcommand{\du}{\mathrm{d}}
\newcommand{\topo}[1]{\mathcal{#1}}

%%% MISCELLANEOUS
\newcommand{\notimplies}{\centernot\implies}

\NewDocumentCommand{\dimu}{O{\mkern-3mu} m}{
\left[
#1
\left[#2\right]
#1
\right]
}

\newcommand{\id}{\operatorname{id}}
\newcommand{\contradiction}{\Rightarrow\!\Leftarrow}
\newcommand{\transition}[3]{{_{#2}\mathbf{#1}_{#3}}}
\newcommand{\barline}{\vspace{10pt}\hrule\vspace{10pt}}
\newcommand{\eqques}{\stackrel{?}{=}}
\newcommand{\Gammaf}[1]{\Gamma\qty(#1)}
\renewcommand\qedsymbol{$\square$}


%%%%%%%%%%%%%%%%%%%%%%%%%%%%%%%%%%%%%%%%%%%%
%       CUSTOM HEADER / SECTION STYLES
%%%%%%%%%%%%%%%%%%%%%%%%%%%%%%%%%%%%%%%%%%%%

\fancypagestyle{titleheader}{
  \fancyhead{}
  \fancyhead{\vspace{2cm}\tline{3pt}\fancyfoot[c]{\thepage}}
  \renewcommand{\headrulewidth}{0pt}
}
\newcommand{\titlebars}{
  \maketitle
  \tline{3pt}
  \thispagestyle{titleheader}
}

%%% CUSTOM SUBSUBSUBSECTION (4-LEVEL HEADINGS)
\titleclass{\subsubsubsection}{straight}[\subsection]
\newcounter{subsubsubsection}[subsubsection]
\renewcommand\thesubsubsubsection{\thesubsubsection.\arabic{subsubsubsection}}
\renewcommand\theparagraph{\thesubsubsubsection.\arabic{paragraph}}
\titleformat{\subsubsubsection}{\normalfont\normalsize\bfseries}{\thesubsubsubsection}{1em}{}
\titlespacing*{\subsubsubsection}{0pt}{3.25ex plus 1ex minus .2ex}{1.5ex plus .2ex}

\makeatletter
\renewcommand\paragraph{\@startsection{paragraph}{5}{\z@}{3.25ex \@plus1ex \@minus.2ex}{-1em}{\normalfont\normalsize\bfseries}}
\renewcommand\subparagraph{\@startsection{subparagraph}{6}{\parindent}{3.25ex \@plus1ex \@minus .2ex}{-1em}{\normalfont\normalsize\bfseries}}
\def\toclevel@subsubsubsection{4}
\def\toclevel@paragraph{5}
\def\toclevel@subparagraph{6}
\def\l@subsubsubsection{\@dottedtocline{4}{7em}{4em}}
\def\l@paragraph{\@dottedtocline{5}{10em}{5em}}
\def\l@subparagraph{\@dottedtocline{6}{14em}{6em}}
\makeatother
\setcounter{secnumdepth}{4}
\setcounter{tocdepth}{4}

%%% CUSTOM SECTION TITLE FONT SIZES
% \titleformat*{\section}{\LARGE\bfseries}
% \titleformat*{\subsection}{\Large\bfseries}
% \titleformat*{\subsubsection}{\large\bfseries}
% \setlength{\parskip}{10pt}

%%%%%%%%%%%%%%%%%%%%%%%%%%%%%%%%%%%%%%%%%%%%
%             ENVIRONMENTS
%%%%%%%%%%%%%%%%%%%%%%%%%%%%%%%%%%%%%%%%%%%%

\newenvironment{amatrix}[1]{%
  \left(
  \begin{array}{@{}*{#1}{c}|c@{}}}{
\end{array}\right)}

\newenvironment{abmatrix}[1]{%
  \left[
  \begin{array}{@{}*{#1}{c}|c@{}}}{
\end{array}\right]}


%%%%%%%%%%%%%%%%%%%%%%%%%%%%%%%%%%%%%%%%%%%%
%          THEOREM STYLES
%%%%%%%%%%%%%%%%%%%%%%%%%%%%%%%%%%%%%%%%%%%%

% \theoremstyle{definition}
% \newtheorem{theorem}{Theorem}[subsection]
% \newtheorem{corollary}{Corollary}[theorem]
% \newtheorem{lemma}[theorem]{Lemma}
% \newtheorem{definition}{Definition}[subsection]
% \newtheorem{example}{Example}[subsection]

% \theoremstyle{remark}
% \newtheorem*{remark}{Remark}
% \newtheorem*{solution}{Solution}
% \newtheorem*{note}{Note}

%%%%%%%%%%%%%%%%%%%%%%%%%%%%%%%%%%%%%%%%%%%%
%        CODE LISTING STYLE
%%%%%%%%%%%%%%%%%%%%%%%%%%%%%%%%%%%%%%%%%%%%


% ============================================
% PYTHON STYLE
% ============================================
\definecolor{vscodeBackground}{RGB}{255,255,255}
\definecolor{vscodeForeground}{RGB}{10,10,10}
\definecolor{vscodeKeyword}{RGB}{60,60,200}
\definecolor{vscodeString}{RGB}{230,0,0}
\definecolor{vscodeComment}{RGB}{0,100,0}
\definecolor{vscodeFunction}{RGB}{220,220,170}
\definecolor{vscodeOperator}{RGB}{212,212,212}
\definecolor{vscodeNumber}{RGB}{181,206,168}

\lstdefinestyle{vscodepython}{
    backgroundcolor=\color{vscodeBackground},
    basicstyle=\ttfamily\footnotesize\color{vscodeForeground},
    commentstyle=\color{vscodeComment},
    keywordstyle=\color{vscodeKeyword}\bfseries,
    stringstyle=\color{vscodeString},
    numberstyle=\tiny\color{vscodeForeground},
    identifierstyle=\color{vscodeForeground},
    morekeywords={True, False, None, and, as, assert, async, await, break, class,
        continue, def, del, elif, else, except, finally, for, from, global, if,
        import, in, is, lambda, nonlocal, not, or, pass, raise, return, try,
        while, with, yield},
    frame=single,
    rulecolor=\color{vscodeForeground},
    breaklines=true,
    numbers=left,
    numbersep=5pt,
    showstringspaces=false,
    tabsize=4,
    captionpos=b,
}

% ============================================
% C/C++ STYLE (FIXED)
% ============================================
\definecolor{vscodeType}{RGB}{0,150,200}
\definecolor{vscodePreprocessor}{RGB}{150,0,150}
\definecolor{vscodeFunction}{RGB}{180,100,0}
\definecolor{vscodeOperator}{RGB}{180,60,180}

\lstdefinestyle{vscodecpp}{
    backgroundcolor=\color{vscodeBackground},
    basicstyle=\ttfamily\scriptsize\color{vscodeForeground},
    commentstyle=\color{vscodeComment},
    stringstyle=\color{vscodeString},
    numberstyle=\tiny\color{vscodeForeground},
    identifierstyle=\color{vscodeForeground},
    %
    % Keywords (class 1)
    keywords={auto, break, case, catch, class, const, constexpr, continue,
        default, delete, do, else, enum, explicit, extern, false, final,
        for, friend, goto, if, inline, mutable, namespace, new, noexcept,
        nullptr, operator, override, private, protected, public, register,
        return, sizeof, static, static_assert, struct, switch, template,
        this, throw, true, try, typedef, typeid, typename, using, virtual,
        volatile, while},
    %
    % Types (class 2) – highlighted with emph
    emph={bool, char, char16_t, char32_t, double, float, int, long, short,
          signed, unsigned, void, wchar_t, size_t, ptrdiff_t},
    emphstyle={\color{vscodeType}\bfseries},
    %
    % Preprocessor (class 3)
    emph={[2]define, elif, else, endif, error, if, ifdef, ifndef, include,
          line, pragma, undef},
    emphstyle={[2]\color{vscodePreprocessor}\bfseries},
    %
    % Standard library functions (class 4)
    emph={[3]printf, scanf, fprintf, fscanf, sprintf, sscanf, getchar, putchar,
          gets, puts, fopen, fclose, fread, fwrite, malloc, calloc, realloc,
          free, exit, abort, assert, system, getenv, qsort, bsearch, memcpy,
          memset, memcmp, strlen, strcpy, strncpy, strcat, strncat, strcmp,
          strncmp, strstr, strtok, atoi, atof, atol, rand, srand, time, clock,
          pow, sqrt, exp, log, log10, sin, cos, tan, asin, acos, atan, atan2,
          fabs, floor, ceil, fmod, feof, ferror, clearerr, remove, rename,
          tmpfile, tmpnam, setbuf, setvbuf, rewind, fgetc, fputc, ungetc,
          fgets, fputs},
    emphstyle={[3]\color{vscodeFunction}\bfseries},
    %
    % C++ extra keywords (add to class 1)
    morekeywords={alignas, alignof, and, and_eq, bitand, bitor, compl,
        concept, const_cast, decltype, dynamic_cast, module, not, not_eq,
        or, or_eq, reinterpret_cast, requires, static_cast, xor, xor_eq},
    %
    % STL types and functions (add to class 4 via emph again)
    emph={[4]std, vector, map, set, string, array, deque, list, queue, stack,
          pair, tuple, unique_ptr, shared_ptr, weak_ptr, make_shared, make_unique,
          move, forward, function, bind, cout, cin, endl, cerr, clog,
          ifstream, ofstream, fstream, stringstream, istringstream, ostringstream,
          sort, find, transform, accumulate, for_each, copy, reverse, unique,
          remove, replace, swap, min, max, lower_bound, upper_bound, binary_search},
    emphstyle={[4]\color{vscodeFunction}\bfseries},
    %
    frame=single,
    rulecolor=\color{vscodeForeground},
    breaklines=true,
    numbers=left,
    numbersep=5pt,
    showstringspaces=false,
    tabsize=4,
    captionpos=b,
}

% ============================================
% SET DEFAULTS
% ============================================
\lstset{style=vscodepython, language=Python}


%% EMPHASIS BOX

\newtcolorbox{emphasis}{
  colback=white,        % White background
  colframe=black,       % Black border
  boxrule=0.5pt,        % Border thickness
  arc=0pt,              % Sharp corners (0pt = no rounding)
  left=6pt,             % Left padding
  right=6pt,            % Right padding
  top=6pt,              % Top padding
  bottom=6pt,           % Bottom padding
  boxsep=0pt,           % No extra separation
  before upper={
    \setlength{\abovedisplayskip}{0pt}
    \setlength{\abovedisplayshortskip}{0pt}
    \setlength{\belowdisplayskip}{0pt}
    \setlength{\belowdisplayshortskip}{0pt}
  }
}

%%%%%%%%%%%%%%%%%%%%%%%%%%%%%%%%%%%%%%%%%%%%
%             DOCUMENT SETTINGS
%%%%%%%%%%%%%%%%%%%%%%%%%%%%%%%%%%%%%%%%%%%%

\setlength{\parindent}{2em}
\setlength{\parskip}{5pt}
\newlength{\originalparindent}
\newlength{\originalparskip}
\AtBeginDocument{%
    \setlength{\originalparindent}{\parindent}%
    \setlength{\originalparskip}{\parskip}%
}
\numberwithin{equation}{section}

% SECTION TITLE FONT SIZES (move here from preset)
\titleformat*{\section}{\LARGE\bfseries}
\titleformat*{\subsection}{\Large\bfseries}
\titleformat*{\subsubsection}{\large\bfseries}
\titleformat*{\subsubsubsection}{\large\bfseries}

% CUSTOM SECTION COMMANDS
\newcommand{\module}[1]{%
    \clearpage
    \refstepcounter{section}
    \vspace*{1cm}
    {\LARGE\bfseries\noindent Module \thesection\par}
    \vspace{0.5cm}
    {\noindent\Huge\bfseries #1\par}
    \vspace{2cm}
    \addcontentsline{toc}{section}{Module \thesection : #1}
}

\newcommand{\unitsec}[1]{%
    \clearpage
    \refstepcounter{section}
    \vspace*{1cm}
    {\LARGE\bfseries\noindent Unit \thesection\par}
    \vspace{0.5cm}
    {\noindent\huge\bfseries #1\par}
    \vspace{1.5cm}
    \addcontentsline{toc}{section}{Unit \thesection : #1}
}

\newcommand{\chaptersec}[2][\relax]{%
    \clearpage
    \refstepcounter{section}
    \vspace*{0.5cm}
    {\LARGE\bfseries\noindent Chapter \thesection\par}
    \vspace{0.5cm}
    {\noindent\huge\bfseries
        \ifx\relax#1\relax
            #2%
        \else
            #1%
        \fi
        \par
    }
    \vspace{1.5cm}
    \addcontentsline{toc}{section}{Chapter \thesection : #2}
}

%%% PROBLEM ENVIRONMENT (was in preset.tex, now here)
\newenvironment{problem}{%
    \hrule
    \vspace{5pt}
}{%
    \vspace{10pt}\hrule\vspace{10pt}
}

\newenvironment{solution_pset}{%
    \noindent\textbf{\Large Solution}\\[0.5cm]
}{%
    \newpage
}


%%% DEFINITION ENVIRONMENT
% Definition counter linked to section
\newcounter{definition}[section]
\renewcommand{\thedefinition}{\thesection.\arabic{definition}}

\newtcolorbox{definition}[2][]{%
    colback=white,
    colframe=black,
    boxrule=0.5pt,
    arc=0pt,
    left=6pt,
    right=6pt,
    top=6pt,
    bottom=6pt,
    fonttitle=\bfseries,
    title={%
        \stepcounter{definition}% Increment counter
        Definition \thedefinition{} : #2% Then display
    },
    #1
}

%%% EXAMPLE ENVIRONMENT
% Example counter linked to section
\newcounter{example}[section]
\renewcommand{\theexample}{\thesection.\arabic{example}}

\newtcolorbox{example}[2][]{%
    breakable,
    colback=white,
    colframe=blue,
    boxrule=0.5pt,
    arc=4pt,
    left=6pt,
    right=6pt,
    top=6pt,
    bottom=6pt,
    fonttitle=\bfseries,
    parbox=false,  % Don't use parbox mode
    before upper={%
        \setlength{\parindent}{\originalparindent}%
        \setlength{\parskip}{\originalparskip}%
        \noindent%
    },
    title={%
        \refstepcounter{example}%
        Example \theexample{} : #2%
    },
    #1
}

%%% THEOREM ENVIRONMENT
% Theorem counter linked to section
\newcounter{theorem}[section]
\renewcommand{\thetheorem}{\thesection.\arabic{theorem}}
\definecolor{darkred}{RGB}{139, 0, 0} 
\newenvironment{theorem}[2][]{%
    \refstepcounter{theorem}%
    \begin{tcolorbox}[
        colback=white,
        colframe=darkred,
        coltitle=white,
        colbacktitle=darkred,
        boxrule=0.5pt,
        arc=0pt,
        left=6pt,
        right=6pt,
        top=6pt,
        bottom=6pt,
        fonttitle=\bfseries,
        title={Theorem \thetheorem{} : #2},
        #1
    ]
}{%
    \end{tcolorbox}%
}

%%% NOTE ENVIRONMENT
\definecolor{darkgray}{RGB}{100, 100, 100}

\newtcolorbox{note}{%
    colback=white,
    colframe=darkgray,
    boxrule=0.5pt,
    arc=0pt,
    left=6pt,
    right=6pt,
    top=6pt,
    bottom=6pt,
    before upper={%
        \noindent{\textit{Note : }}\par\vspace{1em}%
    }
}

%%% REMARK ENVIRONMENT
\newtcolorbox{remark}{%
    colback=white,
    colframe=darkgray,
    boxrule=0.5pt,
    arc=0pt,
    left=6pt,
    right=6pt,
    top=6pt,
    bottom=6pt,
    before upper={%
        \noindent\textbf{\textit{Remark : }}\par\vspace{1em}%
    }
}

\usepackage{bbm}
\begin{document}
	\begin{titlepage}
    \vspace*{2.5in}
    \hrule
    \vspace{3pt}
    \hrule
    \vspace{0.25in}
    \begin{center}
		{
			\Large 
			\textsc{
				University of the Philippines Diliman
			}
			\\[4pt]
			College of Science
			\\[4pt]
			National Institute of Physics
		}
		\\[0.5cm]
		{
			\Huge
			\textbf{
				Physics 180 THV-1
				\\[5pt]
				Problem Set 10
			}
		}
		\\[0.75cm]
		{
			\Large 
			\textsc{
				Vercil Juan
			}
		}\\[5pt] 
		{
			\large 
			IV - BS Physics
		}
    \end{center}
    \vspace{0.25in}
    \hrule
    \vspace{3pt}
    \hrule 
    \vfill
    \begin{center}
    {\large 
		September 24, 2026
	}
    \end{center}
\end{titlepage}

	% ===========================
	% PROBLEM SET 10
	% ===========================
	\stepcounter{section}
	\section*{
		Problem 10.1
	}
	\begin{problem}
		Show that the general solutions of the free-particle Dirac equation for 
		a particle with momentum $\vb p$ are obtainable from the solutions for a 
		particle at rest via Lorentz transformations. Hint: What are the Lorentz 
		transformation properties of the Dirac spinors?

		At rest:
		\begin{equation}
			\label{eq:at-rest}
		    u_1(m,0) 
			=
			\bmqty{
				1\\0\\0\\0
			},\quad  
		    u_2(m,0) 
			=
			\bmqty{
				0\\1\\0\\0
			},\quad  
		    u_3(m,0) 
			=
			\bmqty{
				0\\0\\1\\0
			},\quad  
		    u_4(m,0) 
			=
			\bmqty{
				0\\0\\0\\1
			}
		\end{equation}

		Moving (i.e., there is a boost, like \textit{riding a rocket}):
		\begin{equation}
			\label{eq:moving}
		    u_1
			=
			N
			\bmqty{
				1\\0\\
				\dfrac{p_z}{E+m} 
				\\
				\dfrac{p_x + ip_y}{E+m} 
			},\quad  
		    u_2
			=
			N
			\bmqty{
				0\\1\\
				\dfrac{p_x - ip_y}{E+m} 
				\\
				\dfrac{-p_z}{E+m} 
			},\quad  
		    u_3
			=
			N
			\bmqty{
				\dfrac{p_z}{E+m} 
				\\
				\dfrac{p_x + ip_y}{E+m} 
				\\
				1\\0
			},\quad  
		    u_4
			=
			N
			\bmqty{
				\dfrac{p_x - ip_y}{E+m} 
				\\
				\dfrac{-p_z}{E+m} 
				\\
				0\\1
			}
		\end{equation}
	\end{problem}
	\begin{solution_pset}
		\begin{proof}
			\nocite{flores2026}
			As hinted, a key point is that the Diract spinors 
			transform under the spinor representation of the Lorentz
			group. A boosted solution is therefore obtained by applying 
			the corresponding spinor boost to a rest-frame solultion.

			Consider the four-momentum of a particle at rest, $p_0^\mu$:
			\begin{equation}
				\label{eq:4p-rest}
				p_0^\mu = (m,\vb{0})
			\end{equation}
			
			Boosting it would lead to a new momentum and energy
			\begin{equation}
				\label{eq:4p-moving}
				p^\mu = (E, \vb{p})
			\end{equation}
			where $E = \sqrt{\vb{p}^2 + m^2}$

			The spinor transformation is then 
			\begin{equation}
				\label{eq:spinor-transformation}
				u(p) = S(\Lambda) u(m,0)
			\end{equation}

			Where $S(\Lambda)$ comes from how a Dirac spinor transforms
			\[
				\psi' = S(\Lambda)\psi
			\]
			and has values 
			\[
				S(\Lambda) = \exp(
					- \frac{i}{4}\omega_{\mu\nu}\sigma^{\mu\nu} 
				)
			\]
			where $\omega_{\mu\nu}$ are the six Lorentz parameters, and the spinor generators are:
			\[
				\sigma^{\mu\nu} = \frac{i}{2}[\gamma^\mu,\gamma^\nu] 
			\]

			For a pure boost, the Lorentz parameters are
			$\omega_{0i} = \phi\,\hat{n}_i$ with all $\omega_{ij} = 0$,
			so
			\[
				S(\Lambda)
				=
				\exp\pqty{
					-\frac{i\phi}{2}\,\hat{n}_i\,\sigma^{0i}
				}.
			\]

			Using $\sigma^{0i} = i\gamma^0\gamma^i = i\alpha^i$, where
			$\alpha^i \equiv \gamma^0\gamma^i$, this becomes
			\[
				S(\Lambda)
				=
				\exp\pqty{
					\frac{\phi}{2}\,\vu{n}\cdot\vb*{\alpha}
				}.
			\]

			Since $(\vu{n}\cdot\vb*{\alpha})^2 = I$, the exponential
			linearizes:
			\[
				S(\Lambda)
				=
				\cosh\pqty{\frac{\phi}{2}}I
				+
				\sinh\pqty{\frac{\phi}{2}}
				(\vu{n}\cdot\vb*{\alpha}).
			\]

			At this stage, the transformation is specified by the boost
			direction $\hat{n}$ and rapidity $\phi$. We now trade these for
			the components of the four-momentum $\vb{p}$. The boost takes
			the rest momentum $p_0^\mu = (m,\vb{0})$ to
			$p^\mu = (E,\vb{p})$, so the rapidity satisfies
			\[
				\cosh\phi = \frac{E}{m},
				\qquad
				\sinh\phi = \frac{\abs{\vb{p}}}{m},
			\]
			and the boost direction is $\hat{n} = \vb{p}/\abs{\vb{p}}$. The
			half-angle identities then give
			\[
				\cosh\pqty{\frac{\phi}{2}}
				=
				\sqrt{\frac{E+m}{2m}},
				\qquad
				\sinh\pqty{\frac{\phi}{2}}
				=
				\frac{\abs{\vb{p}}}{\sqrt{2m(E+m)}}.
			\]

			Substituting and factoring out $\sqrt{(E+m)/2m}$, we obtain
			the spinor boost:
			\[
				S(\Lambda)
				=
				\sqrt{\frac{E+m}{2m}}
				\pqty{
					I + \frac{\vb*{\alpha}\cdot\vb{p}}{E+m}
				}.
			\]
			\newpage
			Since a given four-momentum $p^\mu = (E,\vb{p})$ with
			$p^\mu p_\mu = m^2$ is reached from the rest momentum by a
			\emph{unique} pure boost $\Lambda_p$, the transformation is now
			fully determined by $\vb{p}$ alone. We may therefore replace
			the argument $\Lambda$ by the momentum $\vb{p}$ and write the spinor boost matrix as
			\begin{equation}
				\label{eq:Sp}
				S(\vb{p})
				=
				\sqrt{\frac{E+m}{2m}}
				\pqty{
					I + \frac{\vb*{\alpha}\cdot\vb{p}}{E+m}
				}.
			\end{equation}

			In the dirac representation,
			\[
				\vb*{\alpha} = 
				\bmqty{
					0 & \vb*{\sigma}\\
					\vb*{\sigma}& 0
				}, 
			\]
			So, we rewrite the spinor bost matrix as
			\begin{equation}
				\label{eq:Sp2}
				S(\vb{p}) = 
				\sqrt{
					\frac{E+m}{2m} 
				}
				\bmqty{
					I & \frac{\vb*{\sigma}\cdot\vb{p}}{E+m} \\
					\frac{\vb*{\sigma}\cdot\vb{p}}{E+m} & I\\
				} 
			\end{equation}

			Since \cite{flores2026},
			\[
				\vb*{\sigma}\cdot \vb{p}
				=
				\bmqty{
					p_z & p_x - ip_y \\
					p_x + ip_y & -pz
				} 
			\]
			We expand the entire boost operator as a $4\times 4$ matrix
			\[
				S(\vb{p})
				=
				\sqrt{\frac{E+m}{2m}}
				\begin{bmatrix}
					1 & 0 & \dfrac{p_z}{E+m} & \dfrac{p_x - ip_y}{E+m} \\[8pt]
					0 & 1 & \dfrac{p_x + ip_y}{E+m} & \dfrac{-p_z}{E+m} \\[8pt]
					\dfrac{p_z}{E+m} & \dfrac{p_x - ip_y}{E+m} & 1 & 0 \\[8pt]
					\dfrac{p_x + ip_y}{E+m} & \dfrac{-p_z}{E+m} & 0 & 1
				\end{bmatrix}
			\]

			Now we apply this spinor boost operator to the four rest spinors.

			Take $u_1(m,0)$. It is
			\[
				u_1(m,0) = 
				\bmqty{
					1\\0\\0\\0
				} 
			\]

			Applying the spinor boost operator to $u_1$, we obtain
			\begin{align*}
				S(\vb{p}) u_1(m,0)
				&=
				\sqrt{\frac{E+m}{2m}}
				\begin{bmatrix}
					1 & 0 & \dfrac{p_z}{E+m} & \dfrac{p_x - ip_y}{E+m} \\[8pt]
					0 & 1 & \dfrac{p_x + ip_y}{E+m} & \dfrac{-p_z}{E+m} \\[8pt]
					\dfrac{p_z}{E+m} & \dfrac{p_x - ip_y}{E+m} & 1 & 0 \\[8pt]
					\dfrac{p_x + ip_y}{E+m} & \dfrac{-p_z}{E+m} & 0 & 1
				\end{bmatrix}
				\bmqty{
					1\\0\\0\\0
				}\\
				&=
				\sqrt{\frac{E+m}{2m}}
				\bmqty{
					1\\0\\
					\dfrac{p_z}{E+m} 
					\\
					\dfrac{p_x + ip_y}{E+m} 
				}
			\end{align*}

			If we set 
			\[
				N = 
				\sqrt{\frac{E+m}{2m}}
			\]
			Then we obtain the general solution for a boost of $u_1$.
			\begin{equation}
				u_1 = 
				\label{eq:boost-u1}
				N
				\bmqty{
					1\\0\\
					\dfrac{p_z}{E+m} 
					\\
					\dfrac{p_x + ip_y}{E+m} 
				}
			\end{equation}
			Which is precisely what we wanted to show.

			It can be shown trivially\footnote{Which is nerdspeak for ``I am tired of typing this out, and the next answers are obvious'', lmao}.
			that the rest of the spinor basis $u_i, i=2,3,4$ are \textit{unit basis spinors}\footnote{unit basis vector equivalents? Is this what you call it? idk man I'm a relativist, spinors are beyond me. Dunno what those are, but they look like unit basis vectors so yah that's what I'll call em I think} can be written as the $i$-th column 
			of the spinor boost matrix, multiplied by $N$. Hence, it follows that
			\begin{equation}
				\label{eq:sols}
				u_2
				=
				N
				\bmqty{
					0\\1\\
					\dfrac{p_x - ip_y}{E+m} 
					\\
					\dfrac{-p_z}{E+m} 
				},\qquad  
				u_3
				=
				N
				\bmqty{
					\dfrac{p_z}{E+m} 
					\\
					\dfrac{p_x + ip_y}{E+m} 
					\\
					1\\0
				},\qquad  
				u_4
				=
				N
				\bmqty{
					\dfrac{p_x - ip_y}{E+m} 
					\\
					\dfrac{-p_z}{E+m} 
					\\
					0\\1
				}
			\end{equation}
			Thus, we show the general solutions for the free-particle Dirac equation.
		\end{proof}
	\end{solution_pset}
	\newpage

	% ===========================
	\stepcounter{section}
	\bibliographystyle{unsrturl} 
	\bibliography{references} 
	\stepcounter{section}
	% ===========================
	\section*{Appendices}
		\label{appendix}
	
		\subsection*{Code} 
		No code was used for this problem set
	\vfill
	\begin{center}
		\textit{Thank you Mr. Kong}\\
		\vspace{1cm}
		{\large
		\textsc{Physica Vincit Omnia}
		}
	\end{center}
\end{document}
